\documentclass{article}
\usepackage{xspace}
\usepackage{fancyhdr}
\usepackage{graphicx}
\usepackage{amsmath,amssymb,amsthm}
\usepackage{microtype}
\usepackage{enumitem}
\usepackage{ifthen}
\usepackage{totcount}
\usepackage{algpseudocode}
\usepackage[usenames]{xcolor}
\usepackage[most]{tcolorbox}
\usepackage[a4paper, margin=1in]{geometry}

\graphicspath{ {./figures/} }

%--------------------------------------------------------------------------------

\newcommand{\edition}{2026-2027}
\newcommand{\deadline}{27 November 2026, 17:00}
\newcommand{\hwNum}{1}

%--------------------------------------------------------------------------------
% Style stuff
\pagestyle{fancy}
\lhead{Geometric Algorithms (INFOGA \edition)}
\chead{}
\rhead{\mytitle}
\lfoot{INFOGA \edition}
\cfoot{}
\rfoot{\thepage}
\renewcommand{\headrulewidth}{0.4pt}
\renewcommand{\footrulewidth}{0.4pt}
\renewcommand{\familydefault}{\sfdefault}

%--------------------------------------------------------------------------------
% Question macro's

\newcommand{\withpoints}[1]{%
  \addtocounter{pointscounter}{#1} \printpoints{#1}
}
\newcommand{\printpoints}[1]{%
   \ifthenelse{#1 = 0}
              {}
              {\textit{(#1 points)}}\mbox{}
}



\newenvironment{question}[1][0]{\begin{tcolorbox}[parbox=false
                                      ,breakable=true
                                      ,enhanced jigsaw
                                      ,title=\bfseries Question
                                      \stepcounter{questionscounter}\arabic{questionscounter}
                                      \withpoints{#1}]
                         }
                         {\end{tcolorbox}}
\newenvironment{subquestions}{\begin{enumerate}[leftmargin=*
                                        ]}
                             {\end{enumerate}}
\newcommand{\subquestion}[1][0]{\item \withpoints{#1}}

\newtotcounter{questionscounter}
\newtotcounter{pointscounter}

\newcommand{\numquestions}{\total{questionscounter}}
\newcommand{\numpoints}{\total{pointscounter}}

%--------------------------------------------------------------------------------
\newcommand{\myremark}[3]{\textcolor{blue}{\textsc{#1 #2:}} \textcolor{red}{\textsf{#3}}}
% \renewcommand{\myremark}[3]{}
\newcommand{\frank}[2][says]{\myremark{Frank}{#1}{#2}}

%--------------------------------------------------------------------------------
% Theorem Environments
\newtheorem{theorem} {Theorem}
\newtheorem{lemma}[theorem] {Lemma}
\newtheorem{corollary}[theorem] {Corollary}
\newtheorem{problem}[theorem] {Problem}
\newtheorem{observation}[theorem] {Observation}
\newtheorem{claim}[theorem] {Claim}
\newtheorem{invariant}[theorem] {Invariant}

%--------------------------------------------------------------------------------
% Otherwise useful Macro's
\newcommand{\eps}{\ensuremath{\varepsilon}\xspace}
\DeclareMathOperator{\argmin}{argmin}
\DeclareMathOperator{\argmax}{argmax}

\newcommand{\mkmbb}[1]{\ensuremath{\mathbb{#1}}\xspace}
\newcommand{\R}{\mkmbb{R}}

\newcommand{\mathfunc}[1]{\ensuremath{\mathit{#1}}\xspace}

%--------------------------------------------------------------------------------
% The Actual document

\newcommand{\mytitle}{Homework Exam \hwNum\xspace\edition}
\title{\mytitle}
\date{\textbf{Deadline: } \deadline}
\author{{\color{red} My name and StudentID go here!}}


\begin{document}
\maketitle
\thispagestyle{fancy}
  \noindent
  This homework exam has \numquestions\xspace question for a total of
  \numpoints\xspace points. You can earn an additional point by a careful
  preparation of your hand-in: using a good layout, good spelling, good
  figures, no sloppy notation, no statements like ``The algorithm runs in
  $n\log n$.''  (forgetting the $O(..)$ and forgetting to say that it concerns
  time), etc. Use lemmas, theorems, and figures where appropriate.

  \begin{question}[9]
    Let $R$ be a set of $n$ ``red'' points in $\R^2$, and let $B$ be a
    set of $m$ ``blue'' points in $\R^2$. For any two points
    $p,q \in \R^2$ we define their $L_\infty$-distance as
    $d(p,q) = \max\{|p_x-q_x|, |p_y-q_y|\}$.

    \begin{subquestions}

    \item Develop an algorithm that can compute the subset of blue
      points for which there is a red point within $L_\infty$ distance
      one. Your algorithm should run in $O((n+m)\log(n+m))$
      time. Prove that your algorithm is correct and achieves the
      desired running time.

    \item Let $b \in B$ be a blue point, and let $R_b \subseteq R$
      denote the subset of red points that lies within
      $L_\infty$-distance one from $b$. Extend your algorithm above so
      that we can report all such sets $R_b$ in additional $O(k)$
      time, where $k$ is the total size of all sets $R_b$. Prove that
      your algorithm is correct and efficient.

    \item You can actually solve the above problems slightly faster
      than the bounds given above. Argue how / give the slightly
      better bound (and argue why your algorithm achieves it).
    \end{subquestions}
  \end{question}

\end{document}
